% Notes of the Montreal seminar, 1970 (cote 129, typed part): archive cote 129, site page 27 (PDF page 28).
% typescript. Draft compared: none.
% AI-transcribed by Claude Opus 5.5 (claude-opus-5-5, adaptive thinking, effort high), Message Batches API, batch mtl129_full (msgbatch_01AGnDKHwy6ATCM7ehUa6263), 2026-11-13; not yet checked by eye.
\scanpage{129/27}
5.4. \quad Pour tout $n \geq 1$ , on définit les schémas
\[ X^{(p^{n})} = X^{(p^{n}/S)} = [X^{(p^{n-1}/S)}]^{(p/S)} \quad . \]
On note $\mathbb{F}^{n}_{X/S}$ le morphisme composé
\[ \mathbb{F}^{n}_{X/S} : X \xrightarrow{\mathbb{F}_{X/S}} X^{(p)} \xrightarrow{\mathbb{F}_{X^{(p)}/S}} X^{(p^{2})} \to \dots \xrightarrow{\mathbb{F}_{X^{(p^{n-1})}/S}} X^{(p^{n})} \quad . \]
De même, si $G$ est un schéma en groupes commutatifs plat sur $S$ , on
note $\mathbb{V}^{n}_{G}$ le composé
\[ \mathbb{V}^{n}_{G} : G^{(p^{n})} \xrightarrow{\mathbb{V}_{G^{(p^{n-1})}}} G^{(p^{n-1})} \to \dots \xrightarrow{\mathbb{V}_{G}} G \ . \]
On a alors la conséquence facile des relations de 5.3. (correspondant
au cas $n = 1$) :
\[ \mathbb{V}^{n}_{G} \circ \mathbb{F}^{n}_{G/S} = p^{n}\,\mathrm{Id}_{G} \ , \ \mathbb{F}^{n}_{G/S} \circ \mathbb{V}^{n}_{G} = p^{n}\,\mathrm{Id}_{G^{(p^{n})}} \quad . \]
$\mathbb{F}^{n}_{X/S}$ et $\mathbb{V}^{n}_{G}$ s'appellent respectivement le \uline{$n^{\text{ième}}$ itéré du Frobenius re-}%
\uline{latif de} $G$ \uline{sur} $S$ et la \uline{$n^{\text{ième}}$ itéré de la Verschiebung de} $G$ \uline{sur} $S$ .
Dans le cas où la base $S$ est le spectre du corps premier $\mathbb{F}_{p}$ , ce sont
bien les itérés de $\mathbb{F}_{G/S}$ , $\mathbb{V}_{G}$ au sens habituel.

\mspage{\hand{27}}
